Direct Numerical Simulations provide high-fidelity descriptions of
Rayleigh--Taylor instability, but their computational cost limits the number of
available realizations. This constraint motivates the development of generative
surrogate models able to reproduce the statistical properties of simulated
density fields at a lower computational cost. In this work, the objective is to
confront two diffusion-based strategies for the generation of two-dimensional
Rayleigh--Taylor instability fields obtained from DNS: a classical diffusion
model trained directly in physical space, and wavelet-based diffusion models
operating on multi-resolution representations.

The study first introduces the physical and numerical context of the
Rayleigh--Taylor instability and emphasizes its strongly non-linear and
multi-scale nature. Fourier and wavelet representations are then examined as
possible tools for reducing dimensional complexity and capturing scale-dependent
flow structures. The first generative approach is a baseline Denoising
Diffusion Probabilistic Model trained directly in image space on preprocessed
density fields, using physically consistent augmentations based on horizontal
periodicity and symmetry. The second approach introduces wavelet decompositions
in order to explicitly separate coarse structures from localized detail
coefficients before applying diffusion in this transformed space.

The comparison between these approaches is carried out both qualitatively and
quantitatively. Generated ensembles are evaluated using feature-space
distances, including FID and a domain-adapted PhyFID, as well as physically
motivated metrics based on vertical density profiles and fluctuation spectra.
These indicators are designed to assess not only the visual realism of the
generated samples, but also their ability to reproduce statistical and
scale-dependent properties of the reference RTI fields.

The baseline diffusion model generates samples that reproduce the main
organization of the reference simulations and remain significantly closer to
the validation dataset than pure noise or unrelated images. In PhyFID space,
the generated fields obtain a distance of approximately $0.132$ from the
validation set, compared with $6.04$ for pure noise. The physical metrics show
the same trend and confirm that the model captures non-trivial statistical
information about the RTI dataset. The wavelet-based formulation introduces a
physically meaningful multi-resolution inductive bias and shows that diffusion
can learn relationships between large-scale approximations and fine-scale
details. However, in its current implementation, the reconstructed fields remain
rougher than those obtained with the direct image-space diffusion model.

This confrontation therefore provides a first assessment of the benefits and
limitations of wavelet representations for diffusion surrogate modeling of
Rayleigh--Taylor instability. While classical diffusion currently gives the
best generation quality in this study, wavelet-based diffusion remains a
promising direction for improving physical interpretability and scale-aware
generation, provided that the conditioning strategy and architecture are further
refined.
